

\Cref{sec:dynamics} introduces the use of describing functions to quantify the fundamental amplitude of non-sinusoidal force waveforms, which can achieve higher powers than the equivalent sinusoidal waveform for the same force limit.

\citet{hall_open-source_2022} documents the RAFT open-source Python package for offshore wind turbines, which uses the drag describing function with iteration but does not include describing functions for saturation.

\ifdefined\DISSERTATION
    The indirect implementation of describing functions for the force saturation nonlinearity is possible because the describing function quasi-linearized pseudo-spectral method is used for both controller synthesis and evaluation, so the time-domain nonlinear optimal controller that is output from the right of \Cref{fig:mod-freq-domain-synthesis} is fed into the top left of \Cref{fig:mod-freq-domain-evaluation}, resulting in a chaining of the inverse control describing function with the forward control describing function, which cancel out.
    Unless intentionally made different to study robustness, the linearized plant model is identical for the synthesis and evaluation steps.
    In this case, two workflows essentially become one, with no further computation required after the optimal control step.
    Reconstruction of the nonlinear controller is not necessary for the evaluation of the power corresponding to that controller.
    The power can be calculated directly from the frequency-domain response and the describing function of the nonlinear controller.
    The time-domain representation of the nonlinear controller is only necessary for hardware implementation, or if the synthesized controller is to be evaluated with a different method such as time-stepping.
\else
    Because the same describing function is applied in both synthesis and evaluation, the inverse and forward describing functions cancel; the time-domain nonlinear controller need not be reconstructed for power evaluation, only for hardware implementation.
\fi

The simulation permits the nonsinusoidal waveform to have a maximum fundamental a factor of $\frac{4}{\pi}$ above the force limit.
This coincides with the fundamental to peak ratio of a square wave, which is the limiting case both for the bang-singular-bang controller (a piecewise-discontinuous sine-square wave combination that \citet{hendrikx_optimal_2017} finds as the analytically optimal solution), and for the saturated sine controller (a piecewise-continuous sine-square wave combination that \citet{coe_initial_2020} finds as numerically optimal).
In fact, the optimization need not specify the exact nonlinear waveform, since under the describing function filtering hypothesis, only the fundamental matters for simulation purposes.
For hardware implementation purposes, the nonlinear controller can be obtained for the saturated-sine controller by the process outlined in the study \citep{mccabe_force-limited_2024}. 

The appropriate complex scale factor from the optimal unsaturated controller (damping or reactive respectively) is found as described in \Cref{sec:appendix-qp-solution} and will be real for the damping control case and complex in the reactive control case. 

%Mathematically, a saturation factor $0< f_{sat}\leq 1$ is computed as the ratio between the fundamental amplitude of the saturated torque and the original unsaturated torque signal,
%\begin{equation}
%    f_{sat} = \min\left( \alpha\frac{\tau_{max}}{\tau}, 1\right),
%\end{equation}
%where $\tau_{\text{max}}$ is the maximum permissible generator torque and $\alpha$ is the ratio between the fundamental amplitude of the saturated waveform and the maximum torque $\tau_{\text{max}}$, using the describing function derived in \citep{mccabe_force-limited_2024},
%\begin{equation}
%    \alpha = \frac{2}{\pi} \left( \frac{F_p}{\tau_{max}}\sin^{-1}\frac{\tau_{max}}{F_p} +  \sqrt{1 - \left(\frac{\tau_{max}}{F_p}\right)^2} \right).
%\end{equation}
    
%When $\frac{F_p}{\tau_{\text{max}}}$ = 1 (no saturation), then $\alpha= 1$; when $\frac{F_p}{\tau_{\text{max}}} \rightarrow \infty$ (full saturation), then $\alpha= \frac{4}{\pi}$. This makes sense because $\frac{4}{\pi}$ is the harmonic amplitude of a unit square wave, and in full saturation the saturated signal approaches a square wave. $\alpha$ is plotted in \Cref{fig:alpha-fit}.

%The factor $f_{sat}$ describes the ratio of the saturated and unsaturated torque signal. Previous work \citet{mccabe_multidisciplinary_2022} implies that this is also the ratio of the saturated and unsaturated control gains, $m_{sat}$, but this is not the case because the WEC response also changes with saturation. 

% %Substituting $f_{sat}$ back into the equation of motion reveals that the control gain saturation multiplier $m_{sat}$ is the solution to a quadratic equation $a_q m_{sat}^2 + b_q m_{sat} + c_q = 0$. If the uncontrolled dynamics have a natural frequency $\omega_{n,u}$ and damping ratio $\zeta_u$, and the controlled unsaturated dynamics have natural frequency $\omega_n$ and damping ratio $\zeta$, then the quadratic coefficients can be expressed as:

% \begin{equation}\label{eq:abc-matrix}
%     \begin{bmatrix}
%         a_q \\ b_q \\ c_q
%     \end{bmatrix}
%     = 
%     \begingroup % keep the row spacing change local
%     \setlength\arraycolsep{2pt}
%     \begin{bmatrix} -4 & \frac{4}{f_{sat}^2} & \frac{1}{f_{sat}^2} & -1 & 0 & 0\\
%     8 & 0 & 0 & 2 & 8 & 2\\
%     -4 & -4 & -1 & -1 & -8 & -2 \end{bmatrix}
%     \endgroup
%     \begin{bmatrix} 
%  (\theta_1)^2  \\
%   (\theta_2)^2 \\
%    (\theta_3)^2 \\
%    (\theta_4)^2  \\
%    \theta_1 \theta_2 \\  
%     \theta_3 \theta_4
%     \end{bmatrix}
% \end{equation}
% where the basis vector $\vec{\theta}$ is a dimensionless function of the natural frequencies and damping ratios:
% \begin{equation}\label{eq:theta-basis}
%         \begin{bmatrix}
%   \theta_1  \\
%   \theta_2 \\
%    \theta_3 \\
%    \theta_4  \\
%     \end{bmatrix} =     \begin{bmatrix}
% ~ \overline{\omega} (\zeta_u - \overline{\omega}_n\zeta)  \\
%   \overline{\omega} ~\overline{\omega}_n \zeta \\
%   \overline{\omega}^2 - \overline{\omega}_n^2\\
%    \overline{\omega}_n^2 - 1  \\
%     \end{bmatrix}
% \end{equation}
% and the overlined frequencies indicate nondimensionalization by the uncontrolled natural frequency: $\overline{\omega} = \frac{\omega}{\omega_{n,u}}, \overline{\omega}_n = \frac{\omega_n}{\omega_{n,u}}$. The quadratic equation is then solved for $m_{sat}$, and of the two roots, the one that is real-valued and between 0 and 1 is used.

% The no saturation case ($f_{sat} = 1$) for \eqref{abc-matrix} yields $m_{sat}=1$. This confirms the expected behavior that when no force reduction is required, the powertrain coefficients can remain unchanged.

% Equations \eqref{eq:abc-matrix} and \eqref{eq:theta-basis} are an equivalent reformulation of \citet{mccabe_force-limited_2024}, which uses expressions for the complex reflection coefficient $\Gamma$ common in electrical engineering, rather than the damping ratio and natural frequency more common in mechanical engineering. The equivalence is:
% \begin{equation}
%     \Gamma = \frac{z-1}{z+1}, \quad z = \frac{2\theta_1 + i~ \theta_4}{-2(\theta_1+\theta_2) + i~(\theta_3+\theta_4)}
% \end{equation}
% % I am inputting zeta and omega_n corresponding to mechanical dynamcis when calculating m_sat, but I really should be inputting them corresponding to electrical dynamics. Maybe describe that?

% The gain saturation multiplier $m_{sat}$ is then used to modify the powertrain coefficients,
% \begin{align}
%     B_{p,sat} &= m_{sat} B_p, & K_{p,sat} &= m_{sat} K_p.
% \end{align}
% and these new control gains are then plugged into equation \eqref{eq:eom-freq-domain} to find the saturated displacement amplitude $X_{sat}$.
